\section{Relative Frobenius in a composite quantum boundary}
\label{sec:composite-relative-boundary}

Arithmetic composition can retain quantum information that separate
invariant reductions discard. The construction in this section assembles
three actual field registers: two labels and their product. Simultaneous
Frobenius is then removed by invariant projection, while Frobenius on one
factor, transported through multiplication, acts on a relative coordinate.
The relative coordinate is a quantum register of dimension equal to the
extension degree. A specified preparation vanishes at the ordinary
counting endpoint, but its first nonzero grade retains that register and
an experiment detecting both Frobenius and coherence.

The six statements below passed the capped review with a repaired
degree-one scope clause in the spectral addendum. They
concern a conditional quantum boundary of already assembled arithmetic
systems. They impose neither a characteristic-one arithmetic category nor
tensor and direct-sum operations on the resulting collection of boundaries.

\subsection{Arithmetic composition before specialization}

For a named field $E=\mathbb F_{p^r}$, the Hilbert register is
$H_E=\ell^2(E)$ with orthonormal basis $|x\rangle$. Fix a primitive
$p$th root $\zeta_p$ and the additive character
$\psi_E(x)=\zeta_p^{\operatorname{Tr}_{E/\mathbb F_p}(x)}$.
The arithmetic amplitudes include zero preparations, shifts, phases,
Frobenius, the Fourier transform, multiplication into a separate target,
and the named field inclusion and trace-fibre isometries:
\begin{align*}
 X_E(a)|x\rangle&=|x+a\rangle,&
 Z_E(b)|x\rangle&=\psi_E(bx)|x\rangle,\\
 U_E|x\rangle&=|x^p\rangle,&
 F_E|x\rangle&=|E|^{-1/2}\sum_y\psi_E(-xy)|y\rangle,\\
 M_E^{(d)}|x_1,\ldots,x_d,z\rangle
 &=|x_1,\ldots,x_d,z+\prod_jx_j\rangle.
\end{align*}
For a named embedding $i:K\to E$, let
$T_i(y)=i^{-1}(\operatorname{Tr}_{E/i(K)}y)$ and
$\kappa_i=|E|/|K|$. Then
\[
 J_i|a\rangle=|i(a)\rangle,\qquad
 V_i|a\rangle=\kappa_i^{-1/2}\sum_{T_i(y)=a}|y\rangle.
\]
These maps act on ordered independent register words. All square roots
are positive, and the Weyl ordering is $W(a,b)=Z(-b)X(a)$.
\provenance{D8, D1301--D1304, D1306, D1308, FRB-FROB, FRB-TRANSFER}
{foundation.md; transfers-example.md}{arithmetic Frobenius checkers}
{frobenius-hierarchy-adjudication.md}

\begin{definition}[Completed concrete arithmetic processes]
Fix a finite named field diagram in characteristic $p$. Take the actual
matrix image of its arithmetic amplitudes, with scalars in
$\mathbb Q(\zeta_p,\sqrt p)$ and equality of typed actual matrices.
Its finite additive completion has finite lists of objects and rectangular
matrices of arrows. Its self-adjoint projection completion has objects
$(X,e)$ with $e=e^*=e^2$ an actual endomorphism, represented on
$\operatorname{ran}e$, and arrows $f:(X,e)\to(Y,g)$ satisfying $f=gfe$.
Independent tensor uses the tensor of words, projections and arrows.

Retained processes are finite Kraus lists of these arrows with
$\sum_jK_j^*K_j\leq I$; normalized lists have equality. Classical output
tags use block algebras and the sum of ordinary matrix traces. Coherent
additive sums retain the source's rectangular Hom blocks. Additional
reference states are the period references and finite profiles specified
below, together with fixed zero preparations. Collective assembly means
a specified actual isometry, arithmetic circuit and projection.
\end{definition}
\begin{scope}
No surjectivity onto all complex CP maps or new faithfulness theorem for
the older syntactic source is asserted. Coherent sums and classical tags
are different constructions.
\end{scope}
\provenance{D1601}{definitions.md}{composite boundary checker C1, C6}
{composite-boundary-adjudication.md}

\begin{proposition}[Positive processes on completed arithmetic systems]
\statusProved\quad
The completed concrete arithmetic source has finite coherent sums,
independent tensor and self-adjoint retracts. Its Kraus lists realize
completely positive, trace-nonincreasing maps; normalized lists realize
CPTP maps. Their compositions and independent tensors retain these
properties, with every classical outcome kept when specified.
\end{proposition}
\begin{scope}
These are prelimit constructions. Their total descent to the selected
boundary objects is not required or inferred.
\end{scope}
\provenance{CMP-SOURCE}{relational-construction.md, Section 1}
{composite boundary checker C1, C6}{composite-boundary-adjudication.md}

\subsection{A relative coordinate inside a multiplication graph}

\begin{definition}[Relative multiplication-graph register]
Fix a named extension $E/K$ with $|K|=q=p^s$, $|E|=q^n=p^r$ and $r=sn$.
Put $\sigma(x)=x^q$ and $U|x\rangle=|\sigma(x)\rangle$. For every
$\sigma$-orbit $O$ of length $d\mid n$ and $k\in\mathbb Z/d$, put
\[
 v_{O,k}=\frac1{\sqrt d}\sum_{x\in O}
       |x,\sigma^k(x),x\sigma^k(x)\rangle.
\]
Let $G_{d,k}$ be the diagonal projector onto precisely the tuples in these
sums as $O$ varies over length-$d$ orbits. Define
\[
 G=\sum_{d\mid n}\sum_{k=0}^{d-1}G_{d,k},\qquad
 \Delta=U\otimes U\otimes U,\qquad
 T=\frac1n\sum_{j=0}^{n-1}\Delta^j,
\]
and set
\[
 P=TG,\qquad P_d=T\sum_kG_{d,k},\qquad Q_{d,k}=TG_{d,k}.
\]
The relational register is $H_{E/K}^{\rm rel}=\operatorname{ran}P$ with
its inherited inner product. For every $n\geq1$ define the copied graph map
$C:H_E\to H_E^{\otimes3}$ by $C|x\rangle=|x,x,x^2\rangle$.
\end{definition}
\begin{scope}
No orbit origin is chosen. The averages use rational complex scalars;
they never invert $n$ in the finite field. Projector and range properties
are established by the next proposition.
\end{scope}
\provenance{D1602}{definitions.md}{composite boundary checker C2, C3}
{composite-boundary-adjudication.md}

\begin{definition}[Relative Frobenius and common observables]
For the preceding named extension use $M=M_E^{(2)}$ and define
\[
 R=M(I\otimes U\otimes I)M^*,\qquad
 E_{d;a,b}=Q_{d,a}R^{a-b}Q_{d,b}.
\]
Let the common observable algebra be the complex star algebra on
$\operatorname{ran}P$ generated by $R$ and the $Q_{d,k}$. Its proposed
abstract comparison is
\[
 B_n=\bigoplus_{d\mid n}M_d(\mathbb C),\qquad
 e_{d;a,b}\longmapsto E_{d;a,b}.
\]
Write $S_d|k\rangle=|k+1\pmod d\rangle$ and
$u_n^{\rm rel}=\bigoplus_{d\mid n}S_d$. The primitive relative component
is $d=n$; its common algebra is $M_n(\mathbb C)$.
\end{definition}
\begin{scope}
The common algebra does not include arbitrary operators on or between
distinct orbit multiplicities. The arithmetic realization still retains
those multiplicity spaces.
\end{scope}
\provenance{D1603}{definitions.md}{composite boundary checker C2, C3}
{composite-boundary-adjudication.md}

\begin{proposition}[An arithmetic matrix algebra after collective descent]
\statusProved\quad
For every named finite extension $E/K$, the vectors $v_{O,k}$ form an
orthonormal basis of $H_{E/K}^{\rm rel}$. The operators $P_d,Q_{d,k}$
are orthogonal projections. Simultaneous Frobenius fixes this register,
while actual relative Frobenius satisfies
\[
 \Delta v_{O,k}=v_{O,k},\qquad Rv_{O,k}=v_{O,k+1}.
\]
The displayed comparison identifies the common observable algebra with
$B_n$, represented with $c_d(q)/d$ copies of $M_d(\mathbb C)$, where
$c_d(q)=\sum_{a\mid d}\mu(d/a)q^a$ counts exact-period labels.
On every $d>1$ block, $\operatorname{Ad}R$ is a nonidentity CPTP
automorphism and the observables are noncommutative.
\end{proposition}
\begin{scope}
Relative Frobenius is Frobenius on one factor transported through actual
multiplication. It is different from the simultaneous Frobenius removed
by invariant projection.
\end{scope}
\provenance{CMP-REL}{relational-construction.md, Sections 2--3}
{composite boundary checker C2, C3}{composite-boundary-adjudication.md}

\begin{proof}[Arithmetic calculation]
On a tuple, multiplication and its inverse give
\[
 R(x,y,z)=(x,\sigma(y),z-xy+x\sigma(y)).
\]
Thus the multiplication graph is preserved and its relative index advances.
Simultaneous Frobenius permutes the summands of each orbit sum. Its
average is the rank-one projector onto that sum, since the $d$ distinct
tuples repeat $n/d$ times in the average. The graph tests select a relative
index; conjugating them by $R$ supplies every common matrix unit.
In particular $[R,Q_{d,0}]v_{O,0}=v_{O,1}\ne0$ when $d>1$.
\end{proof}

For one orbit, the separately fixed subspace is the line spanned by
$f_O=d^{-1/2}\sum_{x\in O}|x\rangle$. Multiplying two separately fixed
copies into a zero target produces only
\[
 \frac1d\sum_{x,y\in O}|x,y,xy\rangle
       =\frac1{\sqrt d}\sum_kv_{O,k}.
\]
This is one vector in the $d$-dimensional jointly descended register.
The comparison concerns orbitwise invariant reduction, not a claim that
the full single-field observable algebra becomes scalar.

\subsection{A vanishing preparation with a surviving quantum experiment}

For the reference used here, let $h\geq0$, $t=e^h$, and let
$P_{E,a}^{\times}$ project onto nonzero labels of exact absolute Frobenius
period $a\mid r$. Put
\[
 b_1(t)=t-1,\qquad b_a(t)=\sum_{e\mid a}\mu(a/e)t^e\quad(a>1),
 \qquad
 D_E(h)=|0\rangle\langle0|+
      \sum_{a\mid r}\frac{b_a(e^h)}{b_a(p)}P_{E,a}^{\times}.
\]
Then $\rho_E(h)=e^{-rh}D_E(h)$ is the period reference, with
$\rho_E(\log p)=I/|E|$. Its finite profile is
\[
 L_E(h)=|0\rangle\langle0|+hB_E,\qquad
 B_E=\sum_{a\mid r}\frac{\varphi(a)}{b_a(p)}P_{E,a}^{\times},
 \qquad \lambda_E(h)=\frac{L_E(h)}{1+rh}.
\]
For a nonzero positive output series $A(h)=\sum_k h^kA_k$, retain its
first nonzero order $v$, coefficient $M=A_v$, trace $w=\operatorname{Tr}M$
and conditional state $M/w$. For normalizers with constant term one, $w$
is the leading event-probability coefficient. Full profiles are retained
before later conditioning.
\provenance{D1501--D1503, D1506, POS-STATE, POS-FINITE}
{reference-and-protocol.md; finite-operational-boundary.md}
{positive arithmetic checker}{positive-arithmetic-adjudication.md}

\begin{definition}[Copied reference instrument]
For $E/K$ of relative degree $n>1$, let $\Pi_n$ project onto labels of
exact relative Frobenius period $n$. Prepare one $\rho_E(h)$ or
$\lambda_E(h)$ and two fixed zero ancillas. Define the arithmetic isometry
$C|x\rangle=|x,x,x^2\rangle$ by controlled addition and multiplication.
The three preparation amplitudes are
\[
 K_{\rm cut}=I-\Pi_n,\qquad
 K_{\rm fail}=(I-T)C\Pi_n,\qquad K_{\rm ok}=TC\Pi_n.
\]
Their tags are primitive-cut failure in $H_E$, averaging failure in
$H_E^{\otimes3}$, and success in that same three-register word.
After success apply $R$ and measure $Q_{n,1}$ with its complement retained.
In the identity comparison replace $R$ by $I$. In the dephased comparison
insert computational dephasing of all three slots immediately before $R$.
\end{definition}
\begin{scope}
The two fixed zero ancillas are not independently prepared period
references. The final effect includes the invariant projection; a diagonal
graph test alone would not detect the stated coherence.
\end{scope}
\provenance{D1604}{definitions.md}{composite boundary checker C4--C6}
{composite-boundary-adjudication.md}

\begin{definition}[Primitive relative quantum boundary]
Retain the leading event record of the copied reference instrument and
its conditional state on the common $M_n(\mathbb C)$ algebra. The proposed
boundary realization has state $|0\rangle\langle0|$, relative-index
projections, the common matrix units and the channel $\operatorname{Ad}S_n$,
where $S_n$ is the single $n$-cycle. Its arithmetic processes are those
induced by these displayed matrices and specified instruments. A larger
category containing all CP maps is an optional envelope. Physical
multiplicity states at different characteristics are not identified.
\end{definition}
\begin{scope}
No tensor or direct-sum operation between boundary objects is stipulated.
The full physical profile remains necessary for arbitrary later retained
arithmetic continuations.
\end{scope}
\provenance{D1605}{definitions.md}{composite boundary checker C3--C6}
{composite-boundary-adjudication.md}

\begin{theorem}[A quantum Frobenius witness in the first nonzero grade]
\statusProved\quad
For every finite $E/K$ with $|K|=p^s$ and relative degree $n>1$, in every
characteristic, the copied reference instrument is CPTP with its stated
histories. Its successful preparation has probability
\[
 w_{s,n}(h)=\frac{c_n(t^s)}{n t^{sn}},\qquad t=e^h,
 \qquad w_{s,n}(h)=\frac{s\varphi(n)}n h+O(h^2).
\]
For $h>0$, its conditional state on the common $M_n$ block is
$|0\rangle\langle0|$. The final test has conditional probabilities
\[
 \begin{array}{c|ccc}
 \text{circuit}&\text{relative Frobenius}&\text{identity}&\text{dephased}\\\hline
 \text{probability}&1&0&1/n
 \end{array}
\]
The finite-profile preparation has success
$h s\varphi(n)/(n(1+snh))$ and the same leading event records.
\end{theorem}
\begin{scope}
The primitive branch is zero at the ordinary endpoint. The theorem
concerns its specified graded conditional remnant. It preserves the
arithmetic realization and its histories without identifying every
physical state or CP map across characteristics.
\end{scope}
\provenance{CMP-BOUNDARY}{boundary-and-spectrum.md, Sections 1--2}
{composite boundary checker C4--C6}{composite-boundary-adjudication.md}

\begin{proof}[Derivation]
The unnormalized reference mass on the fixed field of $\sigma^d$ is
$t^{sd}$ by reference restriction. Divisor inversion therefore makes
the primitive relative mass $c_n(t^s)$. For $x$ in a primitive orbit,
$TC|x\rangle=v_{O,0}/\sqrt n$, so the diagonal input succeeds with
probability $1/n$ times that mass divided by $t^{sn}$.
Differentiation at zero gives $s\varphi(n)/n$.

Relative Frobenius takes each prepared vector to $v_{O,1}$, while identity
leaves it orthogonal to the final effect. Dephasing replaces each orbit
sum by its uniform computational mixture. Each of its $n$ tuples has
squared overlap $1/n$ with the final invariant vector. The result persists
under the physical mixture over orbits.
\end{proof}

The primitive-cut failure probability is $1-c_n(t^s)/t^{sn}$; averaging
failure has probability $(1-1/n)c_n(t^s)/t^{sn}$. The final test splits
only the success history. Thus every omitted success is accounted for by
a named physical output. The scalar $1/n$ concerns the diagonal reference:
on coherent inputs $K_{\rm ok}^*K_{\rm ok}$ is an orbitwise fixed-vector
projector, not $\Pi_n/n$.

\subsection{The binary quadratic example}

Take $E=\mathbb F_2[\alpha]/(\alpha^2+\alpha+1)$ over $K=\mathbb F_2$.
Its primitive orbit is $\{\alpha,\alpha^2\}$, and the two graph vectors are
\begin{align*}
 v_0&=\frac{|\alpha,\alpha,\alpha^2\rangle+
                  |\alpha^2,\alpha^2,\alpha\rangle}{\sqrt2},\\
 v_1&=\frac{|\alpha,\alpha^2,1\rangle+
                  |\alpha^2,\alpha,1\rangle}{\sqrt2}.
\end{align*}
Their common algebra is a qubit algebra. Actual multiplication-conjugated
relative Frobenius swaps them. The preparation probability is
$w(h)=(1-e^{-h})/2$, hence has coefficient $1/2$ at first order.
At the arithmetic point $t=2$, preparation succeeds with probability
$1/4$. The full final joint success is $1/4$, becomes zero when Frobenius
is removed, and becomes $1/8$ when the orbit coherence is dephased.
These distinctions are generated inside the three arithmetic registers.
\provenance{CMP-REL, CMP-BOUNDARY}{relational-construction.md; boundary-and-spectrum.md}
{composite boundary checker C2--C6}{composite-boundary-adjudication.md}

\subsection{Embedding comparisons and an actual Fourier return}

\begin{definition}[Relational transfers and retained Fourier return]
For a named embedding $i:E\to F$ over the same fixed $K$, restrict
$J_i^{\otimes3}$ to the all-period relational code and call it
$J_i^{\rm rel}$. Its decoder has success amplitude $(J_i^{\rm rel})^*$
and failure amplitude $I-J_i^{\rm rel}(J_i^{\rm rel})^*$, with different
target tags. Towers retain every stopped history and length-$d$ components
remain length $d$ in the upper ambient field.

On a primitive code of degree $n>1$, put $F_{\rm first}=F_E\otimes I\otimes I$.
The return and failure amplitudes are its restrictions followed by $P_n$
and $I-P_n$, respectively; the failure retains the ambient word.
\end{definition}
\begin{scope}
The embedding comparison fixes the base field $K$. No claim about changing
the relative base or Fourier transforms on other slots is included.
\end{scope}
\provenance{D1606}{definitions.md}{composite boundary checker C7, C8}
{composite-boundary-adjudication.md}

\begin{proposition}[Coherent transfers without orbit origins]
\statusProved\quad
Every such named embedding gives an isometry satisfying
$J_i^{\otimes3}v_{O,k}=v_{i(O),k}$. These isometries compose in every
fixed-base tower and intertwine relative Frobenius and the component tests.
Their retained decoders are CPTP, and success maps compose with every
stopped failure history kept distinct.
\end{proposition}
\begin{scope}
A lower primitive component is not replaced by the upper field's
primitive component. No compatible orbit origins are needed, including
for twisted field embeddings.
\end{scope}
\provenance{CMP-NATURAL}{relational-construction.md, Section 4}
{composite boundary checker C7}{composite-boundary-adjudication.md}

\begin{proposition}[First-register Fourier compression]
\statusProved\quad
For every primitive graph code of relative degree $n>1$,
\[
 P_nF_{\rm first}P_n=|E|^{-1/2}D_{\rm phase}P_n,\qquad
 D_{\rm phase}v_{O,k}=\psi_E(-x^2)v_{O,k}\quad(x\in O).
\]
The phase is independent of the selected $x$ in $O$. Retained return has
probability $1/|E|$ for every code state and conditional return is identity
on the common $M_n$ observables.
\end{proposition}
\begin{scope}
The physical multiplicity phases and the unconditioned return probability
retain arithmetic data. The ambient failure branch is not discarded.
\end{scope}
\provenance{CMP-FOURIER}{boundary-and-spectrum.md, Section 3}
{composite boundary checker C8}{composite-boundary-adjudication.md}

\begin{proof}[Kernel calculation]
Fourier in the first slot fixes the other two coordinates $y,z$. On the
primitive graph $y\ne0$ and $z=xy$ forces $x=z/y$. Only the diagonal
Fourier kernel survives compression, namely $|E|^{-1/2}\psi_E(-x^2)$.
Absolute trace is invariant along the relative Frobenius orbit. The
compressed amplitude consequently has squared modulus $1/|E|$ on the
entire code. Complementary projection supplies the other outcome.
\end{proof}

\subsection{Which spectrum and which trace survive}

\begin{definition}[Finite relational spectral data]
Use ordinary Hilbert trace and $\det(I-zR)$ on the all-period relational
code, with actual orbit multiplicities. Keep separately the normalized
trace $\operatorname{tr}_d$ on each common $M_d$ and the linear-map trace
on its channel. To obtain a reference trace on the common blocks, apply
the copied reference and invariant projection, retain success, then
uniformly randomize over $R^j$, $0\leq j<n$, explicitly discarding the
randomization label. The channel can use $n$ copies of $R^j/n$ for each
$j$, so it needs no additional square root in the arithmetic scalar field.
The copied map is defined for every $n\geq1$; moving-sector conditioning
requires $n>1$. This preparation differs from the unrandomized witness.
\end{definition}
\begin{scope}
These definitions concern finite systems. They posit no infinite trace,
regularized determinant, or identification with a Riemann spectral operator.
\end{scope}
\provenance{D1607}{definitions.md}{composite boundary checker C9}
{composite-boundary-adjudication.md}

\begin{proposition}[Finite spectral and reference identities]
\statusProved\quad
For every finite $E/K$ and integer $j$, with $\gcd(n,0)=n$,
\[
 \operatorname{Tr}(R^j)=q^{\gcd(n,j)},\qquad
 \det(I-zR)=\prod_{d\mid n}(1-z^d)^{c_d(q)/d}.
\]
On the full code algebra the channel's linear-map trace is instead
$|\operatorname{Tr}(R^j)|^2$. The primitive implementer has the $n$th
roots of unity as eigenvalues, each once on its common block; its channel
has their ratios, each root with multiplicity $n$.

The randomized copied reference has unnormalized common-block functional
\[
 t^{-sn}\sum_{d\mid n}\frac{c_d(t^s)}d\operatorname{tr}_d(a_d).
\]
For $n>1$, conditioning its moving blocks $d>1$ before taking the first
nonzero grade gives positive normalized weights proportional to $\varphi(d)/d$.
For $n=1$ the moving sector is zero and has no conditional state.
\end{proposition}
\begin{scope}
Ordinary spectral trace and normalized reference expectation are distinct.
The finite determinant identity supplies an arithmetic comparison for
subsequent spectral work; no infinite completion is assumed here.
\end{scope}
\provenance{CMP-TRACE}{boundary-and-spectrum.md, Section 4}
{composite boundary checker C9}{composite-boundary-adjudication.md}

The finite trace identity follows by counting fixed vectors in each
relative cycle. In particular, for $|z|<1$,
\[
 \det(I-zR)^{-1}
 =\exp\!\left(\sum_{j\geq1}\frac{z^j}{j}\operatorname{Tr}(R^j)\right).
\]
The positive quantum remnant and this arithmetic spectral identity are
already properties of an assembled system. Any later tower completion,
rescaling of Frobenius iterates, or new boundary composition must be
specified and justified independently.
